% ============================================================
% Notes: Linear Operators, Self-Adjoint ODEs, Eigenfunctions
% Topics: linear differential operators, Sturm–Liouville,
%         Hermitian/self-adjoint operators, orthogonality,
%         uniqueness, degeneracy
% Source: Arfken & Weber (6th ed.), Ch. 9--10
% ============================================================

\section{Differential equations as linear operators}

\subsection{Linear operators and the superposition principle}
A (differential) operator $L$ is \textbf{linear} if
\begin{equation}
L(a\phi + b\psi) = a\,L(\phi) + b\,L(\psi),
\qquad a,b \in \mathbb{R}\ \text{or}\ \mathbb{C}.
\end{equation}
A linear ODE/PDE can be written abstractly as
\begin{equation}
L\psi = F,
\end{equation}
where $F$ is a known ``source'' function and $\psi$ is the unknown solution.
For the homogeneous equation
\begin{equation}
L\psi = 0,
\end{equation}
the \textbf{superposition principle} holds:
if $\psi_1,\psi_2$ solve $L\psi=0$, then $c_1\psi_1+c_2\psi_2$ also solves it.

\subsection{Second-order linear ODE as an operator equation}
A general second-order linear differential operator can be written as
\begin{equation}
Lu(x) = p_0(x)\,u''(x) + p_1(x)\,u'(x) + p_2(x)\,u(x),
\end{equation}
with real coefficient functions $p_0,p_1,p_2$ (smooth enough on the interval of interest).

\subsection{Eigenvalue problems}
Many physics problems are eigenvalue problems:
\begin{equation}
Lu + \lambda w(x)u = 0,
\qquad w(x)>0\ \text{(except possibly isolated points)}.
\end{equation}
A nontrivial solution $u_\lambda$ that satisfies the required boundary conditions is called an
\textbf{eigenfunction} and the corresponding $\lambda$ is an \textbf{eigenvalue}.
Typically, boundary conditions restrict $\lambda$ to a \textbf{discrete spectrum}.

\paragraph{Quantum mechanics example.}
The time-independent Schr\"odinger equation is an eigenvalue equation:
\begin{equation}
H\psi = E\psi,
\end{equation}
where $H$ is the Hamiltonian operator and $E$ is the energy eigenvalue.

% ------------------------------------------------------------

\section{Self-adjoint ODEs (Sturm--Liouville form)}

\subsection{Adjoint operator and self-adjointness}
Define an inner-product-like bilinear form (real case for simplicity)
\begin{equation}
\langle u,Lu\rangle := \int_a^b u(x)\,Lu(x)\,dx.
\end{equation}
Starting from
\begin{equation}
Lu = p_0u'' + p_1u' + p_2u,
\end{equation}
one defines an \textbf{adjoint operator} $\overline{L}$ such that (schematically)
\begin{equation}
\int_a^b v\,Lu\,dx = \int_a^b (\overline{L}v)\,u\,dx
\quad \text{(plus boundary terms)}.
\end{equation}
If the boundary terms vanish under the imposed boundary conditions and
\begin{equation}
\overline{L} = L,
\end{equation}
then $L$ is called \textbf{self-adjoint}.

\subsection{Self-adjoint (Sturm--Liouville) operator form}
A second-order differential operator is in \textbf{self-adjoint form} if it can be written as
\begin{equation}
Lu = \frac{d}{dx}\!\left(p(x)\frac{du}{dx}\right) + q(x)\,u(x),
\end{equation}
with real functions $p(x),q(x)$.

\subsection{Converting a general ODE to self-adjoint form}
Given
\begin{equation}
Lu = p_0u'' + p_1u' + p_2u,
\end{equation}
one can multiply by a suitable integrating factor to obtain a self-adjoint form.
Define
\begin{equation}
f(x) = \exp\!\left(\int^x \frac{p_1(t)}{p_0(t)}\,dt\right),
\end{equation}
then
\begin{equation}
\frac{f(x)}{p_0(x)}Lu
=
\frac{d}{dx}\!\left(f(x)\frac{du}{dx}\right)
+
\frac{p_2(x)}{p_0(x)}f(x)\,u,
\end{equation}
which is self-adjoint.

\subsection{Sturm--Liouville eigenvalue problem}
The standard Sturm--Liouville problem is
\begin{equation}
\frac{d}{dx}\!\left(p(x)\frac{du}{dx}\right) + q(x)u + \lambda w(x)u = 0,
\qquad a\le x\le b,
\end{equation}
with boundary conditions chosen so that the operator is Hermitian/self-adjoint.

\paragraph{Important note.}
Self-adjointness is not only about the differential expression:
it depends crucially on the \textbf{function space} and \textbf{boundary conditions}.

% ------------------------------------------------------------

\section{Inner product, orthogonality, and normalization}

\subsection{Weighted inner product}
For Sturm--Liouville problems the natural inner product is
\begin{equation}
\langle v|u\rangle := \int_a^b v^*(x)\,w(x)\,u(x)\,dx,
\end{equation}
where $w(x)$ is the weight function.

\subsection{Orthogonality of eigenfunctions}
Let $u_i,u_j$ satisfy
\begin{equation}
Lu_i + \lambda_i w u_i = 0,
\qquad
Lu_j + \lambda_j w u_j = 0.
\end{equation}
If $L$ is Hermitian (with boundary conditions) then for $\lambda_i\neq \lambda_j$:
\begin{equation}
\int_a^b u_i(x)\,u_j^*(x)\,w(x)\,dx = 0.
\end{equation}
So eigenfunctions with distinct eigenvalues are orthogonal under the weight $w$.

\subsection{Normalization}
If $u_n$ is an eigenfunction, scaling is arbitrary for homogeneous equations:
\begin{equation}
u_n \to c\,u_n.
\end{equation}
We often choose orthonormality:
\begin{equation}
\int_a^b u_m(x)\,u_n^*(x)\,w(x)\,dx = \delta_{mn}.
\end{equation}

% ------------------------------------------------------------

\section{Hermitian operators: reality of eigenvalues}

\subsection{Hermitian condition}
An operator $L$ is \textbf{Hermitian} if
\begin{equation}
\int \psi_1^*(x)\,L\psi_2(x)\,dx
=
\int (L\psi_1(x))^*\,\psi_2(x)\,dx
\end{equation}
for all functions in the domain satisfying the boundary conditions.

\subsection{Eigenvalues are real}
If
\begin{equation}
Lu_i + \lambda_i w u_i = 0
\end{equation}
and $L$ is Hermitian, then $\lambda_i\in\mathbb{R}$.
Sketch: multiply the equation by $u_i^*$, subtract the conjugate equation times $u_i$,
integrate, and use the Hermitian property to force
\begin{equation}
\lambda_i^* = \lambda_i.
\end{equation}

\subsection{Expectation values in quantum mechanics}
For a Hermitian operator $L$ and wavefunction $\psi$,
\begin{equation}
\langle L\rangle := \int \psi^*(x)\,L\psi(x)\,dx
\end{equation}
is real, even if $\psi$ is not an eigenfunction.

% ------------------------------------------------------------

\section{Uniqueness, linear independence, and degeneracy}

\subsection{Uniqueness for second-order linear ODEs (idea)}
For a second-order linear homogeneous ODE, the solution space is 2-dimensional:
there are exactly two linearly independent solutions locally (given appropriate regularity).
Imposing two boundary/initial conditions typically selects a unique solution.

\subsection{Linear independence of eigenfunctions (distinct eigenvalues)}
If $L$ has eigenfunctions $u_1,u_2$ with distinct eigenvalues $\lambda_1\neq\lambda_2$,
then $u_1$ and $u_2$ must be linearly independent.
Reason: if $u_2=c\,u_1$, then applying $L$ forces $\lambda_2=\lambda_1$.

\subsection{Degeneracy}
An eigenvalue $\lambda$ is \textbf{$N$-fold degenerate} if there exist $N$ linearly independent
eigenfunctions corresponding to the same eigenvalue.

\paragraph{Key point.}
If $\lambda_i=\lambda_j$ (degenerate case), orthogonality is \emph{not automatic}.
The Sturm--Liouville orthogonality proof only forces orthogonality for
$\lambda_i\neq \lambda_j$.

\subsection{Orthogonalizing degenerate eigenfunctions}
Even if eigenfunctions are degenerate, one can always choose an orthogonal basis
within the degenerate subspace using Gram--Schmidt.

\subsection{Degeneracy from symmetry (physics intuition)}
Degeneracy is often caused by symmetry.
Example: hydrogen atom energy depends on $n$ but not on $(\ell,m)$ in the Coulomb potential,
leading to an $n^2$-fold degeneracy (ignoring spin).
Breaking symmetry (e.g.\ magnetic field) can split levels (Zeeman effect).

% ------------------------------------------------------------

\section{Concrete example: Fourier eigenfunctions (periodic interval)}

Consider
\begin{equation}
\frac{d^2u}{dx^2} + n^2u = 0,
\end{equation}
on an interval of length $2\pi$ with periodic boundary conditions.
Eigenfunctions:
\begin{equation}
u_n(x)=\cos(nx),\qquad v_n(x)=\sin(nx),
\end{equation}
and the eigenvalue is $\lambda=n^2$.
This eigenvalue is \textbf{two-fold degenerate} for $n\ge 1$ because both $\sin(nx)$ and
$\cos(nx)$ correspond to the same $n^2$.

Orthogonality relations (on $[x_0,x_0+2\pi]$):
\begin{align}
\int_{x_0}^{x_0+2\pi}\sin(mx)\sin(nx)\,dx &= C_n\delta_{mn},\\
\int_{x_0}^{x_0+2\pi}\cos(mx)\cos(nx)\,dx &= D_n\delta_{mn},\\
\int_{x_0}^{x_0+2\pi}\sin(mx)\cos(nx)\,dx &= 0.
\end{align}

% ------------------------------------------------------------

\section{What to remember (high-yield summary)}

\begin{itemize}
\item Linear DEs can be written as $L\psi=F$; if $F=0$ then superposition holds.
\item Self-adjoint form is
\[
Lu=(pu')'+qu.
\]
\item Sturm--Liouville eigenvalue equation:
\[
(pu')'+qu+\lambda w u=0.
\]
\item Weighted inner product:
\[
\langle v|u\rangle=\int_a^b v^*wu\,dx.
\]
\item If $L$ is Hermitian/self-adjoint (with boundary conditions):
\begin{itemize}
\item eigenvalues are real,
\item eigenfunctions with distinct eigenvalues are orthogonal,
\item degeneracy means same eigenvalue can have multiple independent eigenfunctions.
\end{itemize}
\item Degenerate eigenfunctions are not automatically orthogonal, but can be orthogonalized.
\end{itemize}
